\heading{数学物理方法（下）-22春-期末}
\noteinfo{题目来源于赛艇，答案由AI生成。}

\begin{question}[20分]

一圆柱体底面半径为 $a$、高为 $h$，其中 $a,h$ 为已知常数。上底面和圆柱侧面均保持零温度，下底面以垂直于底面的恒定热流密度 $q$ 输入热量，其中 $q$ 为已知常数。求达到平衡后圆柱体内温度的稳态分布。

\begin{solution}

取圆柱坐标，稳态温度与角变量无关，记为 $u(\rho,z)$。它满足
\[
\frac1\rho\frac{\partial}{\partial\rho}\!\left(\rho\frac{\partial u}{\partial\rho}\right)+\frac{\partial^2u}{\partial z^2}=0,
\quad 0<\rho<a,\ 0<z<h,
\]
\[
u|_{\rho=a}=0,
\qquad -k\,\frac{\partial u}{\partial z}\Big|_{z=0}=q,
\qquad u|_{z=h}=0.
\]
对 $\rho$ 方向作本征展开，令 $\mu_n$ 为 $J_0(x)$ 的第 $n$ 个正零点，则可得一组本征函数 $J_0(\mu_n\rho/a)$。由边界条件定系数后，稳态解可写为
\ansmath{u(\rho,z)=\frac{2qa}{k}\sum_{n=1}^{\infty}
 \frac{1}{\mu_n^2J_1(\mu_n)}
 \frac{\sinh\!\bigl(\frac{\mu_n}{a}(h-z)\bigr)}{\cosh\!\bigl(\frac{\mu_n}{a}h\bigr)}
 J_0\!\left(\frac{\mu_n}{a}\rho\right)}.
同一个问题也可改写成原题解中的等价形式
\ansmath{u(\rho,z)=\frac{q}{k}(h-z)-\frac{8qh}{\pi^2k}
 \sum_{n=0}^{\infty}
 \frac{1}{(2n+1)^2I_0\!\left(\frac{(2n+1)\pi a}{2h}\right)}
 \cos\!\left(\frac{(2n+1)\pi z}{2h}\right)
 I_0\!\left(\frac{(2n+1)\pi\rho}{2h}\right)}.

\end{solution}
\end{question}

\begin{question}[25分]

求解球内振动问题
\[
\begin{cases}
\dfrac{\partial^2u}{\partial t^2}-a^2\nabla^2u=v_0z, & x^2+y^2+z^2<r_0^2,\ t>0,\\[4pt]
\left.\dfrac{\partial u}{\partial n}\right|_{x^2+y^2+z^2=r_0^2}=0, & t\ge 0,\\[4pt]
u|_{t=0}=0,\qquad \left.\dfrac{\partial u}{\partial t}\right|_{t=0}=0, & x^2+y^2+z^2\le r_0^2,
\end{cases}
\]
其中 $a,v_0,r_0>0$ 为常数。

\begin{solution}

由于右端是 $v_0z=v_0r\cos\theta$，在球坐标下可设
\[
u(r,\theta,t)=v(r,t)\cos\theta.
\]
则径向方程化为
\[
\frac{\partial^2v}{\partial t^2}-a^2\left[\frac1{r^2}\frac{\partial}{\partial r}\left(r^2\frac{\partial v}{\partial r}\right)-\frac{2}{r^2}v\right]=v_0r,
\qquad v_r(r_0,t)=0.
\]
分离变量后得到本征值问题
\[
\frac1{r^2}\frac{d}{dr}\left(r^2\frac{dR}{dr}\right)+\left(\lambda-\frac{2}{r^2}\right)R=0,
\qquad R(0)\text{ 有界},\quad R'(r_0)=0,
\]
其本征函数可取
\[
R_n(r)=j_1\!\left(\frac{\mu_n}{r_0}r\right),
\qquad \lambda_n=\left(\frac{\mu_n}{r_0}\right)^2,
\]
其中 $\mu_n$ 为方程 $j_1'(x)=0$ 的第 $n$ 个正根。把非齐次项按该组本征函数展开，可得
\ansmath{u(r,\theta,t)=\cos\theta\sum_{n=1}^{\infty}
 \frac{2v_0r_0^3j_2(\mu_n)}{\mu_n^2a^2j_1(\mu_n)j_1''(\mu_n)}
 \left(\cos\frac{\mu_nat}{r_0}-1\right)
 j_1\!\left(\frac{\mu_n}{r_0}r\right)}.
利用 $j_1''(\mu_n)=j_1(\mu_n)(2/\mu_n^2-1)$，也可写成
\ansmath{u(r,\theta,t)=\cos\theta\sum_{n=1}^{\infty}
 \frac{2v_0r_0^3j_2(\mu_n)}{\mu_n^2a^2j_1(\mu_n)^2\left(1-2/\mu_n^2\right)}
 \left(1-\cos\frac{\mu_nat}{r_0}\right)
 j_1\!\left(\frac{\mu_n}{r_0}r\right)}.

\end{solution}
\end{question}

\begin{question}[20分]

设一维 Dirac Hamiltonian 算符
\[
\hat H=
\begin{pmatrix}
-i\dfrac{d}{dx} & m_1(x)-im_2(x)\\[4pt]
m_1(x)+im_2(x) & i\dfrac{d}{dx}
\end{pmatrix},
\]
其中 $a\le x\le b$，$m_1(x)$ 和 $m_2(x)$ 都是实值函数。算符作用在两分量自旋子
\[
F(x)=\begin{pmatrix}f_1(x)\\ f_2(x)\end{pmatrix}
\]
上，且 $f_1,f_2\in L^2[a,b]$。内积定义为
\[
\langle F\mid G\rangle=\int_a^b F^{\dagger}(x)G(x)\,dx.
\]
\begin{enumerate}[label=（\arabic*）,leftmargin=2.8em,itemsep=0.2em]
\item 什么条件下算符有伴算符？
\item 分别给出以下两种情形的例子：

情形一：有伴算符但不是自伴算符的 $\hat H$ 及其伴算符 $\hat H^{\dagger}$；

情形二：自伴算符 $\hat H$。
\end{enumerate}

\begin{solution}

对分部积分可得
\[
\langle G\mid \hat HF\rangle-\langle \hat H^{\dagger}G\mid F\rangle
=i\bigl[-g_1^*(x)f_1(x)+g_2^*(x)f_2(x)\bigr]_{x=a}^{x=b}.
\]
因此只要定义域上的边界条件能使该边界项恒为零，算符就存在伴算符。

一个常用充分例子是：若定义域中 $F$ 满足
\ansmath{F(a)=2F(b)},
则伴算符的定义域应取
\ansmath{2G(a)=G(b)}.
于是
\[
\hat H^{\dagger}=
\begin{pmatrix}
-i\dfrac{d}{dx} & m_1(x)-im_2(x)\\[4pt]
m_1(x)+im_2(x) & i\dfrac{d}{dx}
\end{pmatrix},
\qquad D(\hat H)=\{F:F(a)=2F(b)\},
\]
\[
D(\hat H^{\dagger})=\{G:2G(a)=G(b)\},
\]
故这是“有伴但不自伴”的例子。

若改取周期型边界条件
\ansmath{F(a)=F(b)},
则同时需要
\ansmath{G(a)=G(b)},
从而定义域相同，于是
\ansmath{\hat H^{\dagger}=\hat H},
即得到自伴算符的一个例子。

\end{solution}
\end{question}

\begin{question}[20分]

高度为无穷（$-\infty<z<\infty$）的柱体，横截面为如图所示的扇形，其张角为 $\phi_0$、半径为 $a$。柱体一个侧平面绝热，另一个侧平面及圆弧侧面温度恒为 $u_0$，其中 $u_0\ne 0$ 为已知常数。柱体初始温度为
\[
u|_{t=0}=f(x,y,z),
\]
其中 $f$ 为已知实函数。求 $t>0$ 时柱体的温度分布。

\begin{solution}

改用柱坐标 $(\rho,\phi,z)$，并令
\[
u(\rho,\phi,z,t)=v(\rho,\phi,z,t)+u_0.
\]
则 $v$ 满足零边界热方程
\[
\frac{\partial v}{\partial t}-\kappa\left[\frac1\rho\frac\partial{\partial\rho}\left(\rho\frac{\partial v}{\partial\rho}\right)+\frac1{\rho^2}\frac{\partial^2v}{\partial\phi^2}+\frac{\partial^2v}{\partial z^2}\right]=0,
\]
\[
v|_{\phi=0}=0,
\qquad \frac{\partial v}{\partial\phi}\Big|_{\phi=\phi_0}=0,
\qquad v|_{\rho=a}=0,
\qquad v \text{ 在 } z\to\pm\infty \text{ 时有界},
\]
\[
v|_{t=0}=f(\rho,\phi,z)-u_0.
\]
先按角变量展开：
\[
v(\rho,\phi,z,t)=\sum_{n=1}^{\infty}w_n(\rho,z,t)\sin\frac{(2n+1)\pi}{2\phi_0}\phi.
\]
再对 $z$ 作 Fourier 变换，对 $\rho$ 作 Bessel 本征展开。设 $\mu_m^n$ 是
\[
J_{\frac{(2n+1)\pi}{2\phi_0}}(x)=0
\]
的第 $m$ 个正零点，则最终得到
\ansmath{\begin{aligned}
u(\rho,\phi,z,t)=u_0
&+\frac{2}{a^2\phi_0}\sqrt{\frac{2}{\pi\kappa t}}
\sum_{n=1}^{\infty}\sin\frac{(2n+1)\pi}{2\phi_0}\phi
\sum_{m=1}^{\infty}
\frac{e^{-\kappa(\mu_m^n/a)^2t}
J_{\frac{(2n+1)\pi}{2\phi_0}}\!\left(\frac{\mu_m^n}{a}\rho\right)}{\left[J'_{\frac{(2n+1)\pi}{2\phi_0}}(\mu_m^n)\right]^2}\\
&\times \int_0^a\int_{-\infty}^{\infty}\int_0^{\phi_0}
\bigl[f(\rho',\phi',\zeta)-u_0\bigr]
\sin\frac{(2n+1)\pi}{2\phi_0}\phi'
\exp\!\left[-\frac{(z-\zeta)^2}{4\kappa t}\right]
J_{\frac{(2n+1)\pi}{2\phi_0}}\!\left(\frac{\mu_m^n}{a}\rho'\right)
\rho'\,d\phi'\,d\zeta\,d\rho'.
\end{aligned}}
这正是原题解第 11 页给出的紧凑表达式。

\end{solution}
\end{question}

\begin{question}[15分]

求下列定解问题相应的 Green 函数：
\[
\begin{cases}
\nabla^2u+k^2u=f(x,y,z), & -\infty<x<\infty,\ -\infty<y<\infty,\ z>0,\\[4pt]
u|_{z=0}=g(x,y), & -\infty<x<\infty,\ -\infty<y<\infty,\\[4pt]
u \text{ 在无穷远只有会聚波}, & \text{取时间因子为 } y e^{-i\omega t},
\end{cases}
\]
其中 $k>0$ 为已知常数，$f$ 和 $g$ 都是已知函数。

\begin{solution}

设自由空间中满足辐射条件的基本解为
\[
G_1(\mathbf r;\mathbf r')=-\frac{e^{-ik|\mathbf r-\mathbf r'|}}{4\pi|\mathbf r-\mathbf r'|}.
\]
为满足平面 $z=0$ 上的 Dirichlet 边界条件，使用镜像点
\[
\mathbf r'^*=(x',y',-z').
\]
于是取
\[
G(\mathbf r;\mathbf r')=G_1(\mathbf r;\mathbf r')-G_1(\mathbf r;\mathbf r'^*).
\]
写成显式形式就是
\ansmath{G(\mathbf r;\mathbf r')=
\frac{e^{-ik\sqrt{(x-x')^2+(y-y')^2+(z+z')^2}}}{4\pi\sqrt{(x-x')^2+(y-y')^2+(z+z')^2}}
-
\frac{e^{-ik\sqrt{(x-x')^2+(y-y')^2+(z-z')^2}}}{4\pi\sqrt{(x-x')^2+(y-y')^2+(z-z')^2}}}
其中 $-\infty<x,x'<\infty$，$-\infty<y,y'<\infty$，$z,z'>0$。

\end{solution}
\end{question}

